\documentclass[12pt,a4paper]{article}
\usepackage[top=1.8cm,bottom=1.8cm,left=1.8cm,right=1.8cm]{geometry}
\usepackage{fontspec}
\setmainfont{TeX Gyre Termes}
\usepackage{amsmath,amssymb}
\usepackage{fancyhdr}
\usepackage{tasks}
\usepackage{tikz}
\settasks{label=\textbf{\Alph*.},label-width=1.5em,item-indent=2em}
\setlength{\parindent}{0pt}
\setlength{\parskip}{3pt}
\pagestyle{fancy}\fancyhf{}\renewcommand{\headrulewidth}{0pt}
\fancyhead[L]{\small Ngân hàng bài tập Toán 9 -- 9H2-C6-B18-D02}
\cfoot{\small Trang \thepage}
\newif\ifloigiai
\loigiaitrue
\newenvironment{cauhoi}[2]{%
  \par\medskip\noindent\rule{\textwidth}{0.4pt}\par\nopagebreak
  \noindent\textbf{[#1]}\hfill{\scriptsize\ttfamily #2}\par\nopagebreak\smallskip}{\par}
\newcommand{\loigiai}[1]{\ifloigiai\par\smallskip\noindent\textit{Lời giải.} #1\par\fi}
\newcommand{\dapan}[1]{\ifloigiai\par\smallskip\noindent\textit{Đáp án:} \textbf{#1}\par\fi}
\begin{document}
\begin{center}
{\Large\textbf{NGÂN HÀNG BÀI TẬP TOÁN 9}}\\[3pt]
\textbf{Chương 6. Hàm số y = ax² (a ≠ 0). Phương trình bậc hai một ẩn}\\[2pt]
\textbf{Bài 18. Hàm số y = ax² (a ≠ 0)}\\[2pt]
\textbf{Dạng D02: Giao điểm parabol và đường thẳng}\\[2pt]
\textit{Mã dạng: 9H2-C6-B18-D02 \quad 50 bài tự luận + 50 câu trắc nghiệm \quad (TH: 25, VD: 25)}
\end{center}
\textbf{Quy ước mã câu:} \textit{khối, học kì, chương, bài, dạng, loại (TL/TN), mức độ, số thứ tự}.\par Ví dụ: \texttt{9H2-C6-B18-D02-TL-NB-001}.\par Để in bản cho học sinh, đặt \verb|\loigiaifalse| ngay sau \verb|\loigiaitrue|.
\section*{Phần I. Tự luận}
\begin{cauhoi}{Bài 1}{9H2-C6-B18-D02-TL-TH-001}
Tìm toạ độ giao điểm của parabol $(P): y=x^{2}$ và đường thẳng $(d): y=-5x-6$.
\loigiai{Phương trình hoành độ giao điểm: $x^{2}=-5x-6$ hay $x^{2}+5x+6=0$, có hai nghiệm $x=-3$ và $x=-2$. Giao điểm là $(-3;9)$ và $(-2;4)$.}
\end{cauhoi}
\begin{cauhoi}{Bài 2}{9H2-C6-B18-D02-TL-TH-002}
Tìm toạ độ giao điểm của parabol $(P): y=x^{2}$ và đường thẳng $(d): y=1$.
\loigiai{Phương trình hoành độ giao điểm: $x^{2}=0x+1$ hay $x^{2}+0x-1=0$, có hai nghiệm $x=-1$ và $x=1$. Giao điểm là $(-1;1)$ và $(1;1)$.}
\end{cauhoi}
\begin{cauhoi}{Bài 3}{9H2-C6-B18-D02-TL-TH-003}
Tìm toạ độ giao điểm của parabol $(P): y=x^{2}$ và đường thẳng $(d): y=2x+3$.
\loigiai{Phương trình hoành độ giao điểm: $x^{2}=2x+3$ hay $x^{2}-2x-3=0$, có hai nghiệm $x=-1$ và $x=3$. Giao điểm là $(-1;1)$ và $(3;9)$.}
\end{cauhoi}
\begin{cauhoi}{Bài 4}{9H2-C6-B18-D02-TL-TH-004}
Tìm toạ độ giao điểm của parabol $(P): y=x^{2}$ và đường thẳng $(d): y=-x+12$.
\loigiai{Phương trình hoành độ giao điểm: $x^{2}=-x+12$ hay $x^{2}+x-12=0$, có hai nghiệm $x=-4$ và $x=3$. Giao điểm là $(-4;16)$ và $(3;9)$.}
\end{cauhoi}
\begin{cauhoi}{Bài 5}{9H2-C6-B18-D02-TL-TH-005}
Tìm toạ độ giao điểm của parabol $(P): y=x^{2}$ và đường thẳng $(d): y=-x+20$.
\loigiai{Phương trình hoành độ giao điểm: $x^{2}=-x+20$ hay $x^{2}+x-20=0$, có hai nghiệm $x=-5$ và $x=4$. Giao điểm là $(-5;25)$ và $(4;16)$.}
\end{cauhoi}
\begin{cauhoi}{Bài 6}{9H2-C6-B18-D02-TL-TH-006}
Tìm toạ độ giao điểm của parabol $(P): y=x^{2}$ và đường thẳng $(d): y=-6x-5$.
\loigiai{Phương trình hoành độ giao điểm: $x^{2}=-6x-5$ hay $x^{2}+6x+5=0$, có hai nghiệm $x=-5$ và $x=-1$. Giao điểm là $(-5;25)$ và $(-1;1)$.}
\end{cauhoi}
\begin{cauhoi}{Bài 7}{9H2-C6-B18-D02-TL-TH-007}
Tìm toạ độ giao điểm của parabol $(P): y=x^{2}$ và đường thẳng $(d): y=4x-3$.
\loigiai{Phương trình hoành độ giao điểm: $x^{2}=4x-3$ hay $x^{2}-4x+3=0$, có hai nghiệm $x=1$ và $x=3$. Giao điểm là $(1;1)$ và $(3;9)$.}
\end{cauhoi}
\begin{cauhoi}{Bài 8}{9H2-C6-B18-D02-TL-TH-008}
Tìm toạ độ giao điểm của parabol $(P): y=x^{2}$ và đường thẳng $(d): y=5x-6$.
\loigiai{Phương trình hoành độ giao điểm: $x^{2}=5x-6$ hay $x^{2}-5x+6=0$, có hai nghiệm $x=2$ và $x=3$. Giao điểm là $(2;4)$ và $(3;9)$.}
\end{cauhoi}
\begin{cauhoi}{Bài 9}{9H2-C6-B18-D02-TL-TH-009}
Tìm toạ độ giao điểm của parabol $(P): y=x^{2}$ và đường thẳng $(d): y=3x+10$.
\loigiai{Phương trình hoành độ giao điểm: $x^{2}=3x+10$ hay $x^{2}-3x-10=0$, có hai nghiệm $x=-2$ và $x=5$. Giao điểm là $(-2;4)$ và $(5;25)$.}
\end{cauhoi}
\begin{cauhoi}{Bài 10}{9H2-C6-B18-D02-TL-TH-010}
Tìm toạ độ giao điểm của parabol $(P): y=x^{2}$ và đường thẳng $(d): y=-5x-4$.
\loigiai{Phương trình hoành độ giao điểm: $x^{2}=-5x-4$ hay $x^{2}+5x+4=0$, có hai nghiệm $x=-4$ và $x=-1$. Giao điểm là $(-4;16)$ và $(-1;1)$.}
\end{cauhoi}
\begin{cauhoi}{Bài 11}{9H2-C6-B18-D02-TL-TH-011}
Tìm toạ độ giao điểm của parabol $(P): y=x^{2}$ và đường thẳng $(d): y=6x-8$.
\loigiai{Phương trình hoành độ giao điểm: $x^{2}=6x-8$ hay $x^{2}-6x+8=0$, có hai nghiệm $x=2$ và $x=4$. Giao điểm là $(2;4)$ và $(4;16)$.}
\end{cauhoi}
\begin{cauhoi}{Bài 12}{9H2-C6-B18-D02-TL-TH-012}
Tìm toạ độ giao điểm của parabol $(P): y=x^{2}$ và đường thẳng $(d): y=-2x+15$.
\loigiai{Phương trình hoành độ giao điểm: $x^{2}=-2x+15$ hay $x^{2}+2x-15=0$, có hai nghiệm $x=-5$ và $x=3$. Giao điểm là $(-5;25)$ và $(3;9)$.}
\end{cauhoi}
\begin{cauhoi}{Bài 13}{9H2-C6-B18-D02-TL-TH-013}
Tìm toạ độ giao điểm của parabol $(P): y=x^{2}$ và đường thẳng $(d): y=8x-15$.
\loigiai{Phương trình hoành độ giao điểm: $x^{2}=8x-15$ hay $x^{2}-8x+15=0$, có hai nghiệm $x=3$ và $x=5$. Giao điểm là $(3;9)$ và $(5;25)$.}
\end{cauhoi}
\begin{cauhoi}{Bài 14}{9H2-C6-B18-D02-TL-TH-014}
Tìm toạ độ giao điểm của parabol $(P): y=x^{2}$ và đường thẳng $(d): y=-4x$.
\loigiai{Phương trình hoành độ giao điểm: $x^{2}=-4x$ hay $x^{2}+4x+0=0$, có hai nghiệm $x=-4$ và $x=0$. Giao điểm là $(-4;16)$ và $(0;0)$.}
\end{cauhoi}
\begin{cauhoi}{Bài 15}{9H2-C6-B18-D02-TL-TH-015}
Tìm toạ độ giao điểm của parabol $(P): y=x^{2}$ và đường thẳng $(d): y=-3x$.
\loigiai{Phương trình hoành độ giao điểm: $x^{2}=-3x$ hay $x^{2}+3x+0=0$, có hai nghiệm $x=-3$ và $x=0$. Giao điểm là $(-3;9)$ và $(0;0)$.}
\end{cauhoi}
\begin{cauhoi}{Bài 16}{9H2-C6-B18-D02-TL-TH-016}
Tìm toạ độ giao điểm của parabol $(P): y=x^{2}$ và đường thẳng $(d): y=3x$.
\loigiai{Phương trình hoành độ giao điểm: $x^{2}=3x$ hay $x^{2}-3x+0=0$, có hai nghiệm $x=0$ và $x=3$. Giao điểm là $(0;0)$ và $(3;9)$.}
\end{cauhoi}
\begin{cauhoi}{Bài 17}{9H2-C6-B18-D02-TL-TH-017}
Tìm toạ độ giao điểm của parabol $(P): y=x^{2}$ và đường thẳng $(d): y=-6x-8$.
\loigiai{Phương trình hoành độ giao điểm: $x^{2}=-6x-8$ hay $x^{2}+6x+8=0$, có hai nghiệm $x=-4$ và $x=-2$. Giao điểm là $(-4;16)$ và $(-2;4)$.}
\end{cauhoi}
\begin{cauhoi}{Bài 18}{9H2-C6-B18-D02-TL-TH-018}
Tìm toạ độ giao điểm của parabol $(P): y=x^{2}$ và đường thẳng $(d): y=-3x+4$.
\loigiai{Phương trình hoành độ giao điểm: $x^{2}=-3x+4$ hay $x^{2}+3x-4=0$, có hai nghiệm $x=-4$ và $x=1$. Giao điểm là $(-4;16)$ và $(1;1)$.}
\end{cauhoi}
\begin{cauhoi}{Bài 19}{9H2-C6-B18-D02-TL-TH-019}
Tìm toạ độ giao điểm của parabol $(P): y=x^{2}$ và đường thẳng $(d): y=-3x+10$.
\loigiai{Phương trình hoành độ giao điểm: $x^{2}=-3x+10$ hay $x^{2}+3x-10=0$, có hai nghiệm $x=-5$ và $x=2$. Giao điểm là $(-5;25)$ và $(2;4)$.}
\end{cauhoi}
\begin{cauhoi}{Bài 20}{9H2-C6-B18-D02-TL-TH-020}
Tìm toạ độ giao điểm của parabol $(P): y=x^{2}$ và đường thẳng $(d): y=-3x-2$.
\loigiai{Phương trình hoành độ giao điểm: $x^{2}=-3x-2$ hay $x^{2}+3x+2=0$, có hai nghiệm $x=-2$ và $x=-1$. Giao điểm là $(-2;4)$ và $(-1;1)$.}
\end{cauhoi}
\begin{cauhoi}{Bài 21}{9H2-C6-B18-D02-TL-TH-021}
Tìm toạ độ giao điểm của parabol $(P): y=x^{2}$ và đường thẳng $(d): y=x+6$.
\loigiai{Phương trình hoành độ giao điểm: $x^{2}=x+6$ hay $x^{2}-x-6=0$, có hai nghiệm $x=-2$ và $x=3$. Giao điểm là $(-2;4)$ và $(3;9)$.}
\end{cauhoi}
\begin{cauhoi}{Bài 22}{9H2-C6-B18-D02-TL-TH-022}
Tìm toạ độ giao điểm của parabol $(P): y=x^{2}$ và đường thẳng $(d): y=-7x-12$.
\loigiai{Phương trình hoành độ giao điểm: $x^{2}=-7x-12$ hay $x^{2}+7x+12=0$, có hai nghiệm $x=-4$ và $x=-3$. Giao điểm là $(-4;16)$ và $(-3;9)$.}
\end{cauhoi}
\begin{cauhoi}{Bài 23}{9H2-C6-B18-D02-TL-TH-023}
Tìm toạ độ giao điểm của parabol $(P): y=x^{2}$ và đường thẳng $(d): y=-x+2$.
\loigiai{Phương trình hoành độ giao điểm: $x^{2}=-x+2$ hay $x^{2}+x-2=0$, có hai nghiệm $x=-2$ và $x=1$. Giao điểm là $(-2;4)$ và $(1;1)$.}
\end{cauhoi}
\begin{cauhoi}{Bài 24}{9H2-C6-B18-D02-TL-TH-024}
Tìm toạ độ giao điểm của parabol $(P): y=x^{2}$ và đường thẳng $(d): y=2x$.
\loigiai{Phương trình hoành độ giao điểm: $x^{2}=2x$ hay $x^{2}-2x+0=0$, có hai nghiệm $x=0$ và $x=2$. Giao điểm là $(0;0)$ và $(2;4)$.}
\end{cauhoi}
\begin{cauhoi}{Bài 25}{9H2-C6-B18-D02-TL-TH-025}
Tìm toạ độ giao điểm của parabol $(P): y=x^{2}$ và đường thẳng $(d): y=-2x$.
\loigiai{Phương trình hoành độ giao điểm: $x^{2}=-2x$ hay $x^{2}+2x+0=0$, có hai nghiệm $x=-2$ và $x=0$. Giao điểm là $(-2;4)$ và $(0;0)$.}
\end{cauhoi}
\begin{cauhoi}{Bài 26}{9H2-C6-B18-D02-TL-VD-026}
Tìm toạ độ giao điểm của parabol $(P): y=x^{2}$ và đường thẳng $(d): y=-5x$.
\loigiai{Phương trình hoành độ giao điểm: $x^{2}=-5x$ hay $x^{2}+5x+0=0$, có hai nghiệm $x=-5$ và $x=0$. Giao điểm là $(-5;25)$ và $(0;0)$.}
\end{cauhoi}
\begin{cauhoi}{Bài 27}{9H2-C6-B18-D02-TL-VD-027}
Tìm toạ độ giao điểm của parabol $(P): y=x^{2}$ và đường thẳng $(d): y=x+20$.
\loigiai{Phương trình hoành độ giao điểm: $x^{2}=x+20$ hay $x^{2}-x-20=0$, có hai nghiệm $x=-4$ và $x=5$. Giao điểm là $(-4;16)$ và $(5;25)$.}
\end{cauhoi}
\begin{cauhoi}{Bài 28}{9H2-C6-B18-D02-TL-VD-028}
Tìm toạ độ giao điểm của parabol $(P): y=x^{2}$ và đường thẳng $(d): y=7x-10$.
\loigiai{Phương trình hoành độ giao điểm: $x^{2}=7x-10$ hay $x^{2}-7x+10=0$, có hai nghiệm $x=2$ và $x=5$. Giao điểm là $(2;4)$ và $(5;25)$.}
\end{cauhoi}
\begin{cauhoi}{Bài 29}{9H2-C6-B18-D02-TL-VD-029}
Tìm toạ độ giao điểm của parabol $(P): y=x^{2}$ và đường thẳng $(d): y=-7x-10$.
\loigiai{Phương trình hoành độ giao điểm: $x^{2}=-7x-10$ hay $x^{2}+7x+10=0$, có hai nghiệm $x=-5$ và $x=-2$. Giao điểm là $(-5;25)$ và $(-2;4)$.}
\end{cauhoi}
\begin{cauhoi}{Bài 30}{9H2-C6-B18-D02-TL-VD-030}
Tìm toạ độ giao điểm của parabol $(P): y=x^{2}$ và đường thẳng $(d): y=-2x+3$.
\loigiai{Phương trình hoành độ giao điểm: $x^{2}=-2x+3$ hay $x^{2}+2x-3=0$, có hai nghiệm $x=-3$ và $x=1$. Giao điểm là $(-3;9)$ và $(1;1)$.}
\end{cauhoi}
\begin{cauhoi}{Bài 31}{9H2-C6-B18-D02-TL-VD-031}
Tìm toạ độ giao điểm của parabol $(P): y=x^{2}$ và đường thẳng $(d): y=3x-2$.
\loigiai{Phương trình hoành độ giao điểm: $x^{2}=3x-2$ hay $x^{2}-3x+2=0$, có hai nghiệm $x=1$ và $x=2$. Giao điểm là $(1;1)$ và $(2;4)$.}
\end{cauhoi}
\begin{cauhoi}{Bài 32}{9H2-C6-B18-D02-TL-VD-032}
Tìm toạ độ giao điểm của parabol $(P): y=x^{2}$ và đường thẳng $(d): y=-4x+5$.
\loigiai{Phương trình hoành độ giao điểm: $x^{2}=-4x+5$ hay $x^{2}+4x-5=0$, có hai nghiệm $x=-5$ và $x=1$. Giao điểm là $(-5;25)$ và $(1;1)$.}
\end{cauhoi}
\begin{cauhoi}{Bài 33}{9H2-C6-B18-D02-TL-VD-033}
Tìm toạ độ giao điểm của parabol $(P): y=x^{2}$ và đường thẳng $(d): y=16$.
\loigiai{Phương trình hoành độ giao điểm: $x^{2}=0x+16$ hay $x^{2}+0x-16=0$, có hai nghiệm $x=-4$ và $x=4$. Giao điểm là $(-4;16)$ và $(4;16)$.}
\end{cauhoi}
\begin{cauhoi}{Bài 34}{9H2-C6-B18-D02-TL-VD-034}
Tìm toạ độ giao điểm của parabol $(P): y=x^{2}$ và đường thẳng $(d): y=x+2$.
\loigiai{Phương trình hoành độ giao điểm: $x^{2}=x+2$ hay $x^{2}-x-2=0$, có hai nghiệm $x=-1$ và $x=2$. Giao điểm là $(-1;1)$ và $(2;4)$.}
\end{cauhoi}
\begin{cauhoi}{Bài 35}{9H2-C6-B18-D02-TL-VD-035}
Tìm toạ độ giao điểm của parabol $(P): y=x^{2}$ và đường thẳng $(d): y=4$.
\loigiai{Phương trình hoành độ giao điểm: $x^{2}=0x+4$ hay $x^{2}+0x-4=0$, có hai nghiệm $x=-2$ và $x=2$. Giao điểm là $(-2;4)$ và $(2;4)$.}
\end{cauhoi}
\begin{cauhoi}{Bài 36}{9H2-C6-B18-D02-TL-VD-036}
Tìm toạ độ giao điểm của parabol $(P): y=x^{2}$ và đường thẳng $(d): y=-2x+8$.
\loigiai{Phương trình hoành độ giao điểm: $x^{2}=-2x+8$ hay $x^{2}+2x-8=0$, có hai nghiệm $x=-4$ và $x=2$. Giao điểm là $(-4;16)$ và $(2;4)$.}
\end{cauhoi}
\begin{cauhoi}{Bài 37}{9H2-C6-B18-D02-TL-VD-037}
Tìm toạ độ giao điểm của parabol $(P): y=x^{2}$ và đường thẳng $(d): y=5x$.
\loigiai{Phương trình hoành độ giao điểm: $x^{2}=5x$ hay $x^{2}-5x+0=0$, có hai nghiệm $x=0$ và $x=5$. Giao điểm là $(0;0)$ và $(5;25)$.}
\end{cauhoi}
\begin{cauhoi}{Bài 38}{9H2-C6-B18-D02-TL-VD-038}
Tìm toạ độ giao điểm của parabol $(P): y=x^{2}$ và đường thẳng $(d): y=25$.
\loigiai{Phương trình hoành độ giao điểm: $x^{2}=0x+25$ hay $x^{2}+0x-25=0$, có hai nghiệm $x=-5$ và $x=5$. Giao điểm là $(-5;25)$ và $(5;25)$.}
\end{cauhoi}
\begin{cauhoi}{Bài 39}{9H2-C6-B18-D02-TL-VD-039}
Tìm toạ độ giao điểm của parabol $(P): y=x^{2}$ và đường thẳng $(d): y=9$.
\loigiai{Phương trình hoành độ giao điểm: $x^{2}=0x+9$ hay $x^{2}+0x-9=0$, có hai nghiệm $x=-3$ và $x=3$. Giao điểm là $(-3;9)$ và $(3;9)$.}
\end{cauhoi}
\begin{cauhoi}{Bài 40}{9H2-C6-B18-D02-TL-VD-040}
Tìm toạ độ giao điểm của parabol $(P): y=x^{2}$ và đường thẳng $(d): y=3x+4$.
\loigiai{Phương trình hoành độ giao điểm: $x^{2}=3x+4$ hay $x^{2}-3x-4=0$, có hai nghiệm $x=-1$ và $x=4$. Giao điểm là $(-1;1)$ và $(4;16)$.}
\end{cauhoi}
\begin{cauhoi}{Bài 41}{9H2-C6-B18-D02-TL-VD-041}
Tìm toạ độ giao điểm của parabol $(P): y=x^{2}$ và đường thẳng $(d): y=x+12$.
\loigiai{Phương trình hoành độ giao điểm: $x^{2}=x+12$ hay $x^{2}-x-12=0$, có hai nghiệm $x=-3$ và $x=4$. Giao điểm là $(-3;9)$ và $(4;16)$.}
\end{cauhoi}
\begin{cauhoi}{Bài 42}{9H2-C6-B18-D02-TL-VD-042}
Tìm toạ độ giao điểm của parabol $(P): y=x^{2}$ và đường thẳng $(d): y=4x+5$.
\loigiai{Phương trình hoành độ giao điểm: $x^{2}=4x+5$ hay $x^{2}-4x-5=0$, có hai nghiệm $x=-1$ và $x=5$. Giao điểm là $(-1;1)$ và $(5;25)$.}
\end{cauhoi}
\begin{cauhoi}{Bài 43}{9H2-C6-B18-D02-TL-VD-043}
Tìm toạ độ giao điểm của parabol $(P): y=x^{2}$ và đường thẳng $(d): y=4x$.
\loigiai{Phương trình hoành độ giao điểm: $x^{2}=4x$ hay $x^{2}-4x+0=0$, có hai nghiệm $x=0$ và $x=4$. Giao điểm là $(0;0)$ và $(4;16)$.}
\end{cauhoi}
\begin{cauhoi}{Bài 44}{9H2-C6-B18-D02-TL-VD-044}
Tìm toạ độ giao điểm của parabol $(P): y=x^{2}$ và đường thẳng $(d): y=2x+15$.
\loigiai{Phương trình hoành độ giao điểm: $x^{2}=2x+15$ hay $x^{2}-2x-15=0$, có hai nghiệm $x=-3$ và $x=5$. Giao điểm là $(-3;9)$ và $(5;25)$.}
\end{cauhoi}
\begin{cauhoi}{Bài 45}{9H2-C6-B18-D02-TL-VD-045}
Tìm toạ độ giao điểm của parabol $(P): y=x^{2}$ và đường thẳng $(d): y=9x-20$.
\loigiai{Phương trình hoành độ giao điểm: $x^{2}=9x-20$ hay $x^{2}-9x+20=0$, có hai nghiệm $x=4$ và $x=5$. Giao điểm là $(4;16)$ và $(5;25)$.}
\end{cauhoi}
\begin{cauhoi}{Bài 46}{9H2-C6-B18-D02-TL-VD-046}
Tìm toạ độ giao điểm của parabol $(P): y=x^{2}$ và đường thẳng $(d): y=-x+6$.
\loigiai{Phương trình hoành độ giao điểm: $x^{2}=-x+6$ hay $x^{2}+x-6=0$, có hai nghiệm $x=-3$ và $x=2$. Giao điểm là $(-3;9)$ và $(2;4)$.}
\end{cauhoi}
\begin{cauhoi}{Bài 47}{9H2-C6-B18-D02-TL-VD-047}
Tìm toạ độ giao điểm của parabol $(P): y=x^{2}$ và đường thẳng $(d): y=5x-4$.
\loigiai{Phương trình hoành độ giao điểm: $x^{2}=5x-4$ hay $x^{2}-5x+4=0$, có hai nghiệm $x=1$ và $x=4$. Giao điểm là $(1;1)$ và $(4;16)$.}
\end{cauhoi}
\begin{cauhoi}{Bài 48}{9H2-C6-B18-D02-TL-VD-048}
Tìm toạ độ giao điểm của parabol $(P): y=x^{2}$ và đường thẳng $(d): y=2x+8$.
\loigiai{Phương trình hoành độ giao điểm: $x^{2}=2x+8$ hay $x^{2}-2x-8=0$, có hai nghiệm $x=-2$ và $x=4$. Giao điểm là $(-2;4)$ và $(4;16)$.}
\end{cauhoi}
\begin{cauhoi}{Bài 49}{9H2-C6-B18-D02-TL-VD-049}
Tìm toạ độ giao điểm của parabol $(P): y=x^{2}$ và đường thẳng $(d): y=-9x-20$.
\loigiai{Phương trình hoành độ giao điểm: $x^{2}=-9x-20$ hay $x^{2}+9x+20=0$, có hai nghiệm $x=-5$ và $x=-4$. Giao điểm là $(-5;25)$ và $(-4;16)$.}
\end{cauhoi}
\begin{cauhoi}{Bài 50}{9H2-C6-B18-D02-TL-VD-050}
Tìm toạ độ giao điểm của parabol $(P): y=x^{2}$ và đường thẳng $(d): y=-4x-3$.
\loigiai{Phương trình hoành độ giao điểm: $x^{2}=-4x-3$ hay $x^{2}+4x+3=0$, có hai nghiệm $x=-3$ và $x=-1$. Giao điểm là $(-3;9)$ và $(-1;1)$.}
\end{cauhoi}
\clearpage\section*{Phần II. Trắc nghiệm}
\begin{cauhoi}{Câu 1}{9H2-C6-B18-D02-TN-TH-001}
Tìm toạ độ giao điểm của parabol $(P): y=x^{2}$ và đường thẳng $(d): y=-5x-6$.
\begin{tasks}(2)
\task $(3;9),\ (2;4)$
\task $(-3;-9),\ (-2;-4)$
\task $(-3;-3),\ (-2;-2)$
\task $(-3;9),\ (-2;4)$
\end{tasks}
\dapan{D}
\end{cauhoi}
\begin{cauhoi}{Câu 2}{9H2-C6-B18-D02-TN-TH-002}
Tìm toạ độ giao điểm của parabol $(P): y=x^{2}$ và đường thẳng $(d): y=1$.
\begin{tasks}(2)
\task $(-1;1),\ (1;1)$
\task $(1;1),\ (-1;1)$
\task $(-1;-1),\ (1;1)$
\task $(-1;-1),\ (1;-1)$
\end{tasks}
\dapan{A}
\end{cauhoi}
\begin{cauhoi}{Câu 3}{9H2-C6-B18-D02-TN-TH-003}
Tìm toạ độ giao điểm của parabol $(P): y=x^{2}$ và đường thẳng $(d): y=2x+3$.
\begin{tasks}(2)
\task $(-1;-1),\ (3;3)$
\task $(-1;1),\ (3;9)$
\task $(-1;-1),\ (3;-9)$
\task $(1;1),\ (-3;9)$
\end{tasks}
\dapan{B}
\end{cauhoi}
\begin{cauhoi}{Câu 4}{9H2-C6-B18-D02-TN-TH-004}
Tìm toạ độ giao điểm của parabol $(P): y=x^{2}$ và đường thẳng $(d): y=-x+12$.
\begin{tasks}(2)
\task $(-4;-16),\ (3;-9)$
\task $(-4;-4),\ (3;3)$
\task $(4;16),\ (-3;9)$
\task $(-4;16),\ (3;9)$
\end{tasks}
\dapan{D}
\end{cauhoi}
\begin{cauhoi}{Câu 5}{9H2-C6-B18-D02-TN-TH-005}
Tìm toạ độ giao điểm của parabol $(P): y=x^{2}$ và đường thẳng $(d): y=-x+20$.
\begin{tasks}(2)
\task $(-5;-25),\ (4;-16)$
\task $(-5;25),\ (4;16)$
\task $(-5;-5),\ (4;4)$
\task $(5;25),\ (-4;16)$
\end{tasks}
\dapan{B}
\end{cauhoi}
\begin{cauhoi}{Câu 6}{9H2-C6-B18-D02-TN-TH-006}
Tìm toạ độ giao điểm của parabol $(P): y=x^{2}$ và đường thẳng $(d): y=-6x-5$.
\begin{tasks}(2)
\task $(5;25),\ (1;1)$
\task $(-5;-5),\ (-1;-1)$
\task $(-5;25),\ (-1;1)$
\task $(-5;-25),\ (-1;-1)$
\end{tasks}
\dapan{C}
\end{cauhoi}
\begin{cauhoi}{Câu 7}{9H2-C6-B18-D02-TN-TH-007}
Tìm toạ độ giao điểm của parabol $(P): y=x^{2}$ và đường thẳng $(d): y=4x-3$.
\begin{tasks}(2)
\task $(-1;1),\ (-3;9)$
\task $(1;-1),\ (3;-9)$
\task $(1;1),\ (3;9)$
\task $(1;1),\ (3;3)$
\end{tasks}
\dapan{C}
\end{cauhoi}
\begin{cauhoi}{Câu 8}{9H2-C6-B18-D02-TN-TH-008}
Tìm toạ độ giao điểm của parabol $(P): y=x^{2}$ và đường thẳng $(d): y=5x-6$.
\begin{tasks}(2)
\task $(2;-4),\ (3;-9)$
\task $(2;4),\ (3;9)$
\task $(-2;4),\ (-3;9)$
\task $(2;2),\ (3;3)$
\end{tasks}
\dapan{B}
\end{cauhoi}
\begin{cauhoi}{Câu 9}{9H2-C6-B18-D02-TN-TH-009}
Tìm toạ độ giao điểm của parabol $(P): y=x^{2}$ và đường thẳng $(d): y=3x+10$.
\begin{tasks}(2)
\task $(-2;-2),\ (5;5)$
\task $(2;4),\ (-5;25)$
\task $(-2;-4),\ (5;-25)$
\task $(-2;4),\ (5;25)$
\end{tasks}
\dapan{D}
\end{cauhoi}
\begin{cauhoi}{Câu 10}{9H2-C6-B18-D02-TN-TH-010}
Tìm toạ độ giao điểm của parabol $(P): y=x^{2}$ và đường thẳng $(d): y=-5x-4$.
\begin{tasks}(2)
\task $(-4;16),\ (-1;1)$
\task $(-4;-16),\ (-1;-1)$
\task $(-4;-4),\ (-1;-1)$
\task $(4;16),\ (1;1)$
\end{tasks}
\dapan{A}
\end{cauhoi}
\begin{cauhoi}{Câu 11}{9H2-C6-B18-D02-TN-TH-011}
Tìm toạ độ giao điểm của parabol $(P): y=x^{2}$ và đường thẳng $(d): y=6x-8$.
\begin{tasks}(2)
\task $(2;4),\ (4;16)$
\task $(2;-4),\ (4;-16)$
\task $(-2;4),\ (-4;16)$
\task $(2;2),\ (4;4)$
\end{tasks}
\dapan{A}
\end{cauhoi}
\begin{cauhoi}{Câu 12}{9H2-C6-B18-D02-TN-TH-012}
Tìm toạ độ giao điểm của parabol $(P): y=x^{2}$ và đường thẳng $(d): y=-2x+15$.
\begin{tasks}(2)
\task $(-5;25),\ (3;9)$
\task $(-5;-5),\ (3;3)$
\task $(-5;-25),\ (3;-9)$
\task $(5;25),\ (-3;9)$
\end{tasks}
\dapan{A}
\end{cauhoi}
\begin{cauhoi}{Câu 13}{9H2-C6-B18-D02-TN-TH-013}
Tìm toạ độ giao điểm của parabol $(P): y=x^{2}$ và đường thẳng $(d): y=8x-15$.
\begin{tasks}(2)
\task $(3;-9),\ (5;-25)$
\task $(-3;9),\ (-5;25)$
\task $(3;9),\ (5;25)$
\task $(3;3),\ (5;5)$
\end{tasks}
\dapan{C}
\end{cauhoi}
\begin{cauhoi}{Câu 14}{9H2-C6-B18-D02-TN-TH-014}
Tìm toạ độ giao điểm của parabol $(P): y=x^{2}$ và đường thẳng $(d): y=-4x$.
\begin{tasks}(2)
\task $(4;16),\ (0;0)$
\task $(-4;16),\ (0;0)$
\task $(-4;-16),\ (0;0)$
\task $(-4;-4),\ (0;0)$
\end{tasks}
\dapan{B}
\end{cauhoi}
\begin{cauhoi}{Câu 15}{9H2-C6-B18-D02-TN-TH-015}
Tìm toạ độ giao điểm của parabol $(P): y=x^{2}$ và đường thẳng $(d): y=-3x$.
\begin{tasks}(2)
\task $(-3;-3),\ (0;0)$
\task $(-3;9),\ (0;0)$
\task $(3;9),\ (0;0)$
\task $(-3;-9),\ (0;0)$
\end{tasks}
\dapan{B}
\end{cauhoi}
\begin{cauhoi}{Câu 16}{9H2-C6-B18-D02-TN-TH-016}
Tìm toạ độ giao điểm của parabol $(P): y=x^{2}$ và đường thẳng $(d): y=3x$.
\begin{tasks}(2)
\task $(0;0),\ (3;3)$
\task $(0;0),\ (-3;9)$
\task $(0;0),\ (3;9)$
\task $(0;0),\ (3;-9)$
\end{tasks}
\dapan{C}
\end{cauhoi}
\begin{cauhoi}{Câu 17}{9H2-C6-B18-D02-TN-TH-017}
Tìm toạ độ giao điểm của parabol $(P): y=x^{2}$ và đường thẳng $(d): y=-6x-8$.
\begin{tasks}(2)
\task $(-4;-4),\ (-2;-2)$
\task $(-4;16),\ (-2;4)$
\task $(-4;-16),\ (-2;-4)$
\task $(4;16),\ (2;4)$
\end{tasks}
\dapan{B}
\end{cauhoi}
\begin{cauhoi}{Câu 18}{9H2-C6-B18-D02-TN-TH-018}
Tìm toạ độ giao điểm của parabol $(P): y=x^{2}$ và đường thẳng $(d): y=-3x+4$.
\begin{tasks}(2)
\task $(-4;-4),\ (1;1)$
\task $(-4;-16),\ (1;-1)$
\task $(-4;16),\ (1;1)$
\task $(4;16),\ (-1;1)$
\end{tasks}
\dapan{C}
\end{cauhoi}
\begin{cauhoi}{Câu 19}{9H2-C6-B18-D02-TN-TH-019}
Tìm toạ độ giao điểm của parabol $(P): y=x^{2}$ và đường thẳng $(d): y=-3x+10$.
\begin{tasks}(2)
\task $(5;25),\ (-2;4)$
\task $(-5;25),\ (2;4)$
\task $(-5;-5),\ (2;2)$
\task $(-5;-25),\ (2;-4)$
\end{tasks}
\dapan{B}
\end{cauhoi}
\begin{cauhoi}{Câu 20}{9H2-C6-B18-D02-TN-TH-020}
Tìm toạ độ giao điểm của parabol $(P): y=x^{2}$ và đường thẳng $(d): y=-3x-2$.
\begin{tasks}(2)
\task $(-2;4),\ (-1;1)$
\task $(2;4),\ (1;1)$
\task $(-2;-2),\ (-1;-1)$
\task $(-2;-4),\ (-1;-1)$
\end{tasks}
\dapan{A}
\end{cauhoi}
\begin{cauhoi}{Câu 21}{9H2-C6-B18-D02-TN-TH-021}
Tìm toạ độ giao điểm của parabol $(P): y=x^{2}$ và đường thẳng $(d): y=x+6$.
\begin{tasks}(2)
\task $(2;4),\ (-3;9)$
\task $(-2;-2),\ (3;3)$
\task $(-2;-4),\ (3;-9)$
\task $(-2;4),\ (3;9)$
\end{tasks}
\dapan{D}
\end{cauhoi}
\begin{cauhoi}{Câu 22}{9H2-C6-B18-D02-TN-TH-022}
Tìm toạ độ giao điểm của parabol $(P): y=x^{2}$ và đường thẳng $(d): y=-7x-12$.
\begin{tasks}(2)
\task $(-4;-4),\ (-3;-3)$
\task $(-4;-16),\ (-3;-9)$
\task $(4;16),\ (3;9)$
\task $(-4;16),\ (-3;9)$
\end{tasks}
\dapan{D}
\end{cauhoi}
\begin{cauhoi}{Câu 23}{9H2-C6-B18-D02-TN-TH-023}
Tìm toạ độ giao điểm của parabol $(P): y=x^{2}$ và đường thẳng $(d): y=-x+2$.
\begin{tasks}(2)
\task $(-2;4),\ (1;1)$
\task $(-2;-2),\ (1;1)$
\task $(2;4),\ (-1;1)$
\task $(-2;-4),\ (1;-1)$
\end{tasks}
\dapan{A}
\end{cauhoi}
\begin{cauhoi}{Câu 24}{9H2-C6-B18-D02-TN-TH-024}
Tìm toạ độ giao điểm của parabol $(P): y=x^{2}$ và đường thẳng $(d): y=2x$.
\begin{tasks}(2)
\task $(0;0),\ (2;-4)$
\task $(0;0),\ (-2;4)$
\task $(0;0),\ (2;4)$
\task $(0;0),\ (2;2)$
\end{tasks}
\dapan{C}
\end{cauhoi}
\begin{cauhoi}{Câu 25}{9H2-C6-B18-D02-TN-TH-025}
Tìm toạ độ giao điểm của parabol $(P): y=x^{2}$ và đường thẳng $(d): y=-2x$.
\begin{tasks}(2)
\task $(-2;-2),\ (0;0)$
\task $(-2;4),\ (0;0)$
\task $(-2;-4),\ (0;0)$
\task $(2;4),\ (0;0)$
\end{tasks}
\dapan{B}
\end{cauhoi}
\begin{cauhoi}{Câu 26}{9H2-C6-B18-D02-TN-VD-026}
Tìm toạ độ giao điểm của parabol $(P): y=x^{2}$ và đường thẳng $(d): y=-5x$.
\begin{tasks}(2)
\task $(-5;-5),\ (0;0)$
\task $(-5;25),\ (0;0)$
\task $(5;25),\ (0;0)$
\task $(-5;-25),\ (0;0)$
\end{tasks}
\dapan{B}
\end{cauhoi}
\begin{cauhoi}{Câu 27}{9H2-C6-B18-D02-TN-VD-027}
Tìm toạ độ giao điểm của parabol $(P): y=x^{2}$ và đường thẳng $(d): y=x+20$.
\begin{tasks}(2)
\task $(-4;-4),\ (5;5)$
\task $(-4;16),\ (5;25)$
\task $(4;16),\ (-5;25)$
\task $(-4;-16),\ (5;-25)$
\end{tasks}
\dapan{B}
\end{cauhoi}
\begin{cauhoi}{Câu 28}{9H2-C6-B18-D02-TN-VD-028}
Tìm toạ độ giao điểm của parabol $(P): y=x^{2}$ và đường thẳng $(d): y=7x-10$.
\begin{tasks}(2)
\task $(2;2),\ (5;5)$
\task $(2;-4),\ (5;-25)$
\task $(-2;4),\ (-5;25)$
\task $(2;4),\ (5;25)$
\end{tasks}
\dapan{D}
\end{cauhoi}
\begin{cauhoi}{Câu 29}{9H2-C6-B18-D02-TN-VD-029}
Tìm toạ độ giao điểm của parabol $(P): y=x^{2}$ và đường thẳng $(d): y=-7x-10$.
\begin{tasks}(2)
\task $(-5;25),\ (-2;4)$
\task $(5;25),\ (2;4)$
\task $(-5;-25),\ (-2;-4)$
\task $(-5;-5),\ (-2;-2)$
\end{tasks}
\dapan{A}
\end{cauhoi}
\begin{cauhoi}{Câu 30}{9H2-C6-B18-D02-TN-VD-030}
Tìm toạ độ giao điểm của parabol $(P): y=x^{2}$ và đường thẳng $(d): y=-2x+3$.
\begin{tasks}(2)
\task $(3;9),\ (-1;1)$
\task $(-3;9),\ (1;1)$
\task $(-3;-9),\ (1;-1)$
\task $(-3;-3),\ (1;1)$
\end{tasks}
\dapan{B}
\end{cauhoi}
\begin{cauhoi}{Câu 31}{9H2-C6-B18-D02-TN-VD-031}
Tìm toạ độ giao điểm của parabol $(P): y=x^{2}$ và đường thẳng $(d): y=3x-2$.
\begin{tasks}(2)
\task $(-1;1),\ (-2;4)$
\task $(1;-1),\ (2;-4)$
\task $(1;1),\ (2;2)$
\task $(1;1),\ (2;4)$
\end{tasks}
\dapan{D}
\end{cauhoi}
\begin{cauhoi}{Câu 32}{9H2-C6-B18-D02-TN-VD-032}
Tìm toạ độ giao điểm của parabol $(P): y=x^{2}$ và đường thẳng $(d): y=-4x+5$.
\begin{tasks}(2)
\task $(-5;25),\ (1;1)$
\task $(5;25),\ (-1;1)$
\task $(-5;-25),\ (1;-1)$
\task $(-5;-5),\ (1;1)$
\end{tasks}
\dapan{A}
\end{cauhoi}
\begin{cauhoi}{Câu 33}{9H2-C6-B18-D02-TN-VD-033}
Tìm toạ độ giao điểm của parabol $(P): y=x^{2}$ và đường thẳng $(d): y=16$.
\begin{tasks}(2)
\task $(-4;16),\ (4;16)$
\task $(4;16),\ (-4;16)$
\task $(-4;-16),\ (4;-16)$
\task $(-4;-4),\ (4;4)$
\end{tasks}
\dapan{A}
\end{cauhoi}
\begin{cauhoi}{Câu 34}{9H2-C6-B18-D02-TN-VD-034}
Tìm toạ độ giao điểm của parabol $(P): y=x^{2}$ và đường thẳng $(d): y=x+2$.
\begin{tasks}(2)
\task $(-1;-1),\ (2;2)$
\task $(1;1),\ (-2;4)$
\task $(-1;-1),\ (2;-4)$
\task $(-1;1),\ (2;4)$
\end{tasks}
\dapan{D}
\end{cauhoi}
\begin{cauhoi}{Câu 35}{9H2-C6-B18-D02-TN-VD-035}
Tìm toạ độ giao điểm của parabol $(P): y=x^{2}$ và đường thẳng $(d): y=4$.
\begin{tasks}(2)
\task $(-2;-4),\ (2;-4)$
\task $(-2;-2),\ (2;2)$
\task $(-2;4),\ (2;4)$
\task $(2;4),\ (-2;4)$
\end{tasks}
\dapan{C}
\end{cauhoi}
\begin{cauhoi}{Câu 36}{9H2-C6-B18-D02-TN-VD-036}
Tìm toạ độ giao điểm của parabol $(P): y=x^{2}$ và đường thẳng $(d): y=-2x+8$.
\begin{tasks}(2)
\task $(-4;-16),\ (2;-4)$
\task $(4;16),\ (-2;4)$
\task $(-4;-4),\ (2;2)$
\task $(-4;16),\ (2;4)$
\end{tasks}
\dapan{D}
\end{cauhoi}
\begin{cauhoi}{Câu 37}{9H2-C6-B18-D02-TN-VD-037}
Tìm toạ độ giao điểm của parabol $(P): y=x^{2}$ và đường thẳng $(d): y=5x$.
\begin{tasks}(2)
\task $(0;0),\ (-5;25)$
\task $(0;0),\ (5;5)$
\task $(0;0),\ (5;25)$
\task $(0;0),\ (5;-25)$
\end{tasks}
\dapan{C}
\end{cauhoi}
\begin{cauhoi}{Câu 38}{9H2-C6-B18-D02-TN-VD-038}
Tìm toạ độ giao điểm của parabol $(P): y=x^{2}$ và đường thẳng $(d): y=25$.
\begin{tasks}(2)
\task $(-5;-5),\ (5;5)$
\task $(-5;25),\ (5;25)$
\task $(-5;-25),\ (5;-25)$
\task $(5;25),\ (-5;25)$
\end{tasks}
\dapan{B}
\end{cauhoi}
\begin{cauhoi}{Câu 39}{9H2-C6-B18-D02-TN-VD-039}
Tìm toạ độ giao điểm của parabol $(P): y=x^{2}$ và đường thẳng $(d): y=9$.
\begin{tasks}(2)
\task $(-3;9),\ (3;9)$
\task $(3;9),\ (-3;9)$
\task $(-3;-3),\ (3;3)$
\task $(-3;-9),\ (3;-9)$
\end{tasks}
\dapan{A}
\end{cauhoi}
\begin{cauhoi}{Câu 40}{9H2-C6-B18-D02-TN-VD-040}
Tìm toạ độ giao điểm của parabol $(P): y=x^{2}$ và đường thẳng $(d): y=3x+4$.
\begin{tasks}(2)
\task $(-1;-1),\ (4;-16)$
\task $(-1;-1),\ (4;4)$
\task $(1;1),\ (-4;16)$
\task $(-1;1),\ (4;16)$
\end{tasks}
\dapan{D}
\end{cauhoi}
\begin{cauhoi}{Câu 41}{9H2-C6-B18-D02-TN-VD-041}
Tìm toạ độ giao điểm của parabol $(P): y=x^{2}$ và đường thẳng $(d): y=x+12$.
\begin{tasks}(2)
\task $(3;9),\ (-4;16)$
\task $(-3;-9),\ (4;-16)$
\task $(-3;-3),\ (4;4)$
\task $(-3;9),\ (4;16)$
\end{tasks}
\dapan{D}
\end{cauhoi}
\begin{cauhoi}{Câu 42}{9H2-C6-B18-D02-TN-VD-042}
Tìm toạ độ giao điểm của parabol $(P): y=x^{2}$ và đường thẳng $(d): y=4x+5$.
\begin{tasks}(2)
\task $(-1;-1),\ (5;5)$
\task $(-1;-1),\ (5;-25)$
\task $(1;1),\ (-5;25)$
\task $(-1;1),\ (5;25)$
\end{tasks}
\dapan{D}
\end{cauhoi}
\begin{cauhoi}{Câu 43}{9H2-C6-B18-D02-TN-VD-043}
Tìm toạ độ giao điểm của parabol $(P): y=x^{2}$ và đường thẳng $(d): y=4x$.
\begin{tasks}(2)
\task $(0;0),\ (4;16)$
\task $(0;0),\ (4;4)$
\task $(0;0),\ (4;-16)$
\task $(0;0),\ (-4;16)$
\end{tasks}
\dapan{A}
\end{cauhoi}
\begin{cauhoi}{Câu 44}{9H2-C6-B18-D02-TN-VD-044}
Tìm toạ độ giao điểm của parabol $(P): y=x^{2}$ và đường thẳng $(d): y=2x+15$.
\begin{tasks}(2)
\task $(-3;-9),\ (5;-25)$
\task $(-3;-3),\ (5;5)$
\task $(-3;9),\ (5;25)$
\task $(3;9),\ (-5;25)$
\end{tasks}
\dapan{C}
\end{cauhoi}
\begin{cauhoi}{Câu 45}{9H2-C6-B18-D02-TN-VD-045}
Tìm toạ độ giao điểm của parabol $(P): y=x^{2}$ và đường thẳng $(d): y=9x-20$.
\begin{tasks}(2)
\task $(-4;16),\ (-5;25)$
\task $(4;16),\ (5;25)$
\task $(4;4),\ (5;5)$
\task $(4;-16),\ (5;-25)$
\end{tasks}
\dapan{B}
\end{cauhoi}
\begin{cauhoi}{Câu 46}{9H2-C6-B18-D02-TN-VD-046}
Tìm toạ độ giao điểm của parabol $(P): y=x^{2}$ và đường thẳng $(d): y=-x+6$.
\begin{tasks}(2)
\task $(-3;-9),\ (2;-4)$
\task $(-3;9),\ (2;4)$
\task $(3;9),\ (-2;4)$
\task $(-3;-3),\ (2;2)$
\end{tasks}
\dapan{B}
\end{cauhoi}
\begin{cauhoi}{Câu 47}{9H2-C6-B18-D02-TN-VD-047}
Tìm toạ độ giao điểm của parabol $(P): y=x^{2}$ và đường thẳng $(d): y=5x-4$.
\begin{tasks}(2)
\task $(-1;1),\ (-4;16)$
\task $(1;1),\ (4;4)$
\task $(1;1),\ (4;16)$
\task $(1;-1),\ (4;-16)$
\end{tasks}
\dapan{C}
\end{cauhoi}
\begin{cauhoi}{Câu 48}{9H2-C6-B18-D02-TN-VD-048}
Tìm toạ độ giao điểm của parabol $(P): y=x^{2}$ và đường thẳng $(d): y=2x+8$.
\begin{tasks}(2)
\task $(-2;-4),\ (4;-16)$
\task $(-2;-2),\ (4;4)$
\task $(2;4),\ (-4;16)$
\task $(-2;4),\ (4;16)$
\end{tasks}
\dapan{D}
\end{cauhoi}
\begin{cauhoi}{Câu 49}{9H2-C6-B18-D02-TN-VD-049}
Tìm toạ độ giao điểm của parabol $(P): y=x^{2}$ và đường thẳng $(d): y=-9x-20$.
\begin{tasks}(2)
\task $(5;25),\ (4;16)$
\task $(-5;-5),\ (-4;-4)$
\task $(-5;25),\ (-4;16)$
\task $(-5;-25),\ (-4;-16)$
\end{tasks}
\dapan{C}
\end{cauhoi}
\begin{cauhoi}{Câu 50}{9H2-C6-B18-D02-TN-VD-050}
Tìm toạ độ giao điểm của parabol $(P): y=x^{2}$ và đường thẳng $(d): y=-4x-3$.
\begin{tasks}(2)
\task $(-3;-9),\ (-1;-1)$
\task $(-3;9),\ (-1;1)$
\task $(3;9),\ (1;1)$
\task $(-3;-3),\ (-1;-1)$
\end{tasks}
\dapan{B}
\end{cauhoi}
\clearpage\section*{Bảng đáp án trắc nghiệm}
\begin{center}\begin{tabular}{|c|c|c|c|c|c|c|c|c|c|}
\hline
1. D & 2. A & 3. B & 4. D & 5. B & 6. C & 7. C & 8. B & 9. D & 10. A\\\hline
11. A & 12. A & 13. C & 14. B & 15. B & 16. C & 17. B & 18. C & 19. B & 20. A\\\hline
21. D & 22. D & 23. A & 24. C & 25. B & 26. B & 27. B & 28. D & 29. A & 30. B\\\hline
31. D & 32. A & 33. A & 34. D & 35. C & 36. D & 37. C & 38. B & 39. A & 40. D\\\hline
41. D & 42. D & 43. A & 44. C & 45. B & 46. B & 47. C & 48. D & 49. C & 50. B\\\hline
\end{tabular}\end{center}
\end{document}
